\chapter{Antimitochondrial Antibody (AMA)}

\section*{What Is This Test?}
The antimitochondrial antibody test detects autoantibodies directed against mitochondrial components. The most clinically important subtype is AMA-M2, which is strongly associated with primary biliary cholangitis (PBC), an autoimmune disease affecting the small bile ducts in the liver. AMA is interpreted with liver enzymes, especially alkaline phosphatase, symptoms, and other investigations.

\section*{Alternative Names}
\begin{itemize}
  \item AMA
  \item AMA-M2
  \item Antimitochondrial Antibody Test
  \item M2 Antigen Test
  \item Anti-Mitochondrial Antibody
  \item PBC Antibody Test
\end{itemize}
\section*{Why It Is Ordered}
Evaluation of unexplained or persistent cholestatic liver-test abnormalities, particularly an elevated alkaline phosphatase; investigation of itching, fatigue, jaundice, or other features suggesting PBC; and clarification of autoimmune liver disease when viral, obstructive, or metabolic causes are being considered.

\section*{How This Test Is Performed}
Blood is collected by venipuncture. AMA may be detected by indirect immunofluorescence and/or measured with an immunoassay that identifies the M2 antigen. The laboratory reports the result as negative, positive, or with a titer or concentration according to the assay.

\section*{Preparation, Risks, and Limitations}
Fasting is usually not required, although the laboratory or clinician may give different instructions. Blood collection may cause brief pain, bruising, or lightheadedness. A positive AMA strongly supports PBC in the appropriate clinical setting but does not by itself establish disease severity or replace clinical evaluation. Rare positive results occur in other autoimmune or liver conditions, and a negative result does not exclude every form of autoimmune cholestatic disease.

\section*{References}
\begin{enumerate}
  \item MedlinePlus. Antimitochondrial Antibody.
  \item European Association for the Study of the Liver. Clinical practice guidelines for primary biliary cholangitis.
\end{enumerate}

\clearpage
\section*{Understanding Results and Follow-up}
The laboratory report should identify the specimen, method, units, reference interval or decision limit, and whether the result is positive, negative, abnormal, or inconclusive. Reference intervals vary with age, sex, pregnancy, specimen type, assay, and laboratory; a result outside a range is not automatically a diagnosis, and a result within a range does not always exclude disease. Discuss unexpected results with the ordering clinician, who can decide whether repeat or confirmatory testing, treatment, or referral is appropriate. Do not change medicines or delay urgent care based on an isolated result.


\section*{Regulatory Status and Authorization Date}
This chapter describes a test category, not a specific commercial product. FDA, CDSCO, EU IVDR, and MHRA approval or registration dates therefore cannot be assigned without the assay manufacturer, product name, instrument, and intended use. For laboratory-developed tests, an individual agency approval date may not exist. See the Preface for the interpretation of regulatory status.

\clearpage
